\begin{tikzpicture}[font=\small,entity/.style={draw=celadon!65,fill=pale!25,text width=3.2cm,align=left,inner sep=7pt}]
\node[entity](course) at (0,2.1){\textbf{课程与知识}\\课程ID、知识ID\\章节、先修边};
\node[entity](question) at (4.9,2.1){\textbf{资源与题目}\\所属知识、版本\\答案、量规、状态};
\node[entity](attempt) at (9.8,2.1){\textbf{学生作答}\\题目ID、任务ID\\得分、提示、时间};
\node[entity](task) at (0,-.6){\textbf{教师任务}\\课程、目标学生\\发布时基线};
\node[entity](evidence) at (4.9,-.6){\textbf{证据更新}\\作答ID、权重\\前后状态、更新依据};
\node[entity](state) at (9.8,-.6){\textbf{知识与复习状态}\\学生ID、知识ID\\参数、间隔与时间};
\draw(course)--node[above]{1:N}(question);\draw(question)--node[above]{1:N}(attempt);
\draw(course)--node[left]{1:N}(task);\draw(task)--node[above]{关联}(evidence);
\draw(attempt)--(evidence);\draw(evidence)--node[above]{聚合}(state);\draw(attempt)--node[right]{更新}(state);
\end{tikzpicture}
